\documentclass{article}
\usepackage{amsfonts}
\usepackage{amsmath}
\usepackage{amssymb}
\newcommand{\sz}{\spadesuit}
\newcommand{\hz}{\heartsuit}
\newcommand{\dz}{\diamondsuit}
\newcommand{\cz}{\clubsuit}
\newcommand{\st}{\rule{0pt}{4mm}}

\begin{document}
\title{The Polish $1\cz$ Opening \\ `Manchester Style'}
\author{These notes by David Collier \\ System created by David Collier and Mike Bell}
\maketitle

\tableofcontents

\newpage
\section{Introduction}
\subsection{A Polish Club system}

The opening bids in our system will look like this:

\begin{center}
\begin{tabular}{|l|p{10cm}|} \hline
$1\cz$ & \st Three-way:
\begin{itemize}
\item 12-14 HCP, balanced or 4=4=1=4 or precisely 5 clubs.
\item 15+ HCP with 5+ clubs.
\item 18+ HCP, any distribution (but some hands open $1\diamondsuit$ or $1\hz$ instead).
\end{itemize} \\ \hline
$1\dz$ & \st 11-21 HCP with 5+ diamonds, or 4-4-4-1 with 4 diamonds, or 4 diamonds and 5 clubs. \\ \hline
$1\hz$ & \st 11-21 HCP with 5+ hearts. \\ \hline
$1\sz$ & \st 11-17 HCP with 5+ spades. \\ \hline
1NT & \st (14)15-17 HCP, frequently off-shape (including 4=4=1=4). \\ \hline
$2\cz$ & \st 10-14 HCP with 6+ clubs, or 5 clubs and a 4-card major. \\ \hline
\end{tabular}
\end{center}

\vspace{3mm}
There are two distinguishing features of this style (though these are far from original). Firstly, balanced hands in the 12-14 HCP range are always opened $1\cz$ unless we have a 5-card suit that we want to bid instead. This is different to how it is done in WJ05, the most widely-known Polish Club variant, which opens $1\dz$ on balanced hands with 4 diamonds. And secondly, not all hands of 18+ HCP are opened $1\cz$. Unless the hand is worth a game force, we would prefer to open $1\dz$ on an unbalanced hand with primary diamonds. Similarly, when the longest suit is hearts we would prefer to open $1\hz$ if the hand is not single-suited. So while the $1\sz$ opening is limited to 17 HCP, the $1\dz$ and $1\hz$ openings are not.

Opening bids of $2\dz$ and higher are pre-emptive; the system does not specify which scheme of pre-empts should be used.

\subsection{About these notes}

These notes are not a complete system, since we will only be looking at auctions after the $1\cz$ opening: continutations after the other opening bids are not defined. This is because our $1\hz$, $1\sz$ and 1NT openings in particular are very similar to `standard' openings playing strong NT and 5-card majors, and so any normal continuations will work.\footnote{Of course, there will be some slight changes that need to be made. For example, because the $1\sz$ opening is limited to 17 HCP, the standard meaning of $1\sz$:1NT,3x does not make sense -- this can be played instead as showing a good two-suiter, 5-5 or better.} For example after a major-suit opening, some people like to play a two-over-one response as game forcing, whereas others prefer lighter responses. We don't want to make life difficult for people wanting to play Polish Club by imposing new responses to familiar opening bids. Similarly, we will not be looking at slam-bidding conventions since these are fairly independent of the the basic structure over $1\cz$.

This is a system for serious partnerships. The aim is to provide efficient agreements for every situation that you are likely to come across, including detailed notes on competitive auctions. There will be many situations where we give precise meanings to calls that might be poorly defined in natural methods, and we will also be using many fairly complex artificial sequences in places where natural bidding would be inefficient. Examples of unusual sequences in this system include:
\begin{itemize}
\item Artificial continuations after opener's $2\cz$ and $2\dz$ rebids. (Opener's $2\dz$ here is usually an artificial game force.)
\item Artificial ways for responder to show minor-suit hands in response to $1\cz$, with artificial rebids based on showing shortage.
\item After an overcall, a mixture of natural bids and transfers by responder.
\item The use of opener's $3\dz$ rebid in competition to show certain strong hands (often a strong hand with clubs, if $3\cz$ would be non-forcing).
\item Some specialised Lebensohl and scrambling 2NT sequences.
\end{itemize}

\vspace{4mm}
The usual note on abbreviations: `M' refers to an unspecified major suit; if a major has already been bid, then `M' normally means `the first-bid major' while `OM' is `the other major'. Similarly we have `m' and `om' for minor suits. When describing hand patterns, `3=4=1=5' means a hand with precisely 3 spades, 4 hearts, 1 diamond and 5 clubs. To give a hand pattern without specifying the order of the suits, I'll put hyphens between the numbers as in `4-4-4-1'.


\newpage
\section{Choice of opening bid}

As explained in the outline of the system, the $1\cz$ opening is a three-way bid. The three options are:\footnote{When describing the bid to opponents, it seems to be easier to explain the natural option as starting at 12 HCP (but denying 6 clubs if minimum), so that you do not have to mention the possibility of a club suit twice. However, the hands of 12-14 HCP with 5 clubs are generally bid as if they were `weak NT' hands in our system, so for our purposes these hands will always be thought of as belonging to the weak type.}

\begin{itemize}
\item The `weak' type -- 12-14 HCP, balanced / 4=4=1=4 / precisely 5 clubs.
\item The `natural' type -- 15+ HCP with 5+ clubs.
\item The `strong' type -- Various 18+ HCP hands.
\end{itemize}

However the decision as to whether to open $1\cz$ is rather more complicated than can be explained with a simple definition, and we need to look at each of the three possibilities in more detail.

\subsection{The `weak' hands}

The stated point range is 12-14, but `good' 14-point hands can be upgraded to a 1NT opening. At the lower end of the range, upgrading 11-counts to $1\cz$ is rather less common, and there are some 12-counts that should be passed: it is important to keep the $1\cz$ opening up to strength. So while we may be playing $1\dz$, $1\hz$ and $1\sz$ as `light' openings, the $1\cz$ opening should be a fairly strict 12 HCP.

One of the most important features of our style of Polish Club is that \emph{all} 4-3-3-3 and 4-4-3-2 hands in the `weak NT' range are opened $1\cz$ -- even hands with a weak doubleton club such as $\sz$KJx $\hz$QJxx $\dz$AQJx $\cz$xx.\footnote{In third and fourth seat, you might agree that it is acceptable to open a good 4-card suit on a balanced hand. But this possibility would not affect how responder bids over $1\cz$.} This is, of course, very different to the more common `better minor' style of playing 5-card majors. For us, the weak option of $1\cz$ simply implies a balanced (or close-to-balanced) hand; it does not suggest clubs over any other suit.

With 5-3-3-2 shape, we have a choice between opening $1\cz$ (to show the balanced hand) and opening the 5-card suit. When the suit is a major it would be normal to open 1M, but you might make an exception if the suit is poor. For 5-3-3-2 hands with 5 diamonds it is a more difficult decision. In a natural-based system it would be clear to open $1\dz$, but playing Polish Club $1\cz$ is made more attractive because you are showing the balanced nature of the hand immediately. So while $\sz$Axx $\hz$xx $\dz$KJxxx $\cz$Axx should still be opened $1\dz$, a hand like $\sz$KJx $\hz$AQ $\dz$Qxxxx $\cz$JTx looks more like a $1\cz$ bid. These decisions may also depend on the continuations you use after a natural opening bid, which might present you with a rebid problem on certain 5-3-3-2 shapes.

The weak option of the $1\cz$ opening also contains some hands which are not normally thought of as balanced. $1\cz$ is the systemic bid for hands in the weak NT range with 4=4=1=4 shape, and also for hands with precisely 5 clubs not suitable for another suit opening. But when you open $1\cz$ with one of these hands, the plan is to bid it as if it was a `weak NT'. In some auctions you may be able to reveal the unbalanced nature of the hand, but much of the time -- and particularly in competitive auctions -- this will not be possible. In particular you need to be aware that when you open $1\cz$ with five clubs on a hand in the weak NT range, you may not be able to show the club suit later. You have to stick with the decision to bid the hand as if it was balanced.

On hands with 5 clubs and a 5-card major, we always open the major. Similarly, with 5 clubs and 4 diamonds the system bid is $1\dz$, but there is some room for judgement here and it is reasonable to open $1\cz$ instead when the diamond suit is particularly weak, for example with $\sz$KJ $\hz$Kx $\dz$xxxx $\cz$AQxxx.

Hands with 5 clubs and a 4-card major pose the most interesting problems, and what happens here will depend on how you like to define your $2\cz$ opening bid. It is perfectly possible to define the $2\cz$ opening as promising at least 6 clubs, putting all the hands with 5 clubs and a 4-card major into $1\cz$. The more traditional approach is include 5$\cz$-4M hands in the $2\cz$ opening, but even if you play this style there will be some hands where the club suit is too weak to open $2\cz$, and so the $1\cz$ opening is needed for those. The precise shape of the hand will also make a difference, and we should look at the possible types individually.

\begin{description}
\item[2=4=2=5 or 4=2=2=5] These semi-balanced hands are very good for opening $1\cz$, since they deliver the two cards in every suit that partner will be expecting. At the same time, these hands are not very good for opening $2\cz$ because of the relatively poor playing strength. So, even if you have agreed to open $2\cz$ on 5$\cz$-4M hands, a $1\cz$ opening would be more normal with these semi-balanced shapes. You would open $2\cz$ only if the club suit was especially good, or if the values were concentrated in the long suits, for example $\sz$AKJx $\hz$xx $\dz$xx $\cz$KQTxx.

As was said earlier, it is important to keep the $1\cz$ opening up to strength, and so despite the fact that a 5-4-2-2 shape is worth more than a typical balanced hand, we would still pass most 11-counts. And if we have a hand which is too weak to open $1\cz$, it should not be opened $2\cz$ either: while a $2\cz$ opening is allowed on fewer then 12 HCP, this would almost always require a 6-card or longer club suit except in third seat.

\item[3=4=1=5, 4=3=1=5 or 4=4=0=5] Despite the shortage, these hands are still relatively good for opening $1\cz$. Partner will be aware of the possibility of shortage in diamonds (remember that 4=4=1=4 hands are also opened $1\cz$), and so he will not bid a diamond suit quite as freely as he would bid clubs or a major. Also, while we do not like having to treat these hands as balanced, we will rarely actually have to rebid in no-trumps because we can raise any major-suit reponse. And these hands are not particularly good for opening $2\cz$ either, because with length in both majors there is a significant chance of missing a major-suit fit. The way to find your fits in the majors is to open $1\cz$. So we will choose $1\cz$ nearly as often as for the semi-balanced shapes.

The 4=4=0=5 hand is particularly difficult no matter how you open it, but if you open $1\cz$ you do give yourself the option of upgrading it to the natural option if partner shows a major.

\item[4=1=3=5 or 1=4=3=5] These are the worst hands for the system. They are not good for opening $1\cz$ because if partner has a major suit, he will be expecting more than a singleton in support: you may end up playing in a 5-1 fit. And with a 1=4=3=5 hand you would have to rebid 1NT on a singleton after $1\cz$ : $1\sz$. So, if a $2\cz$ opening is allowed on these hands, this will be the normal choice unless the club suit is particularly poor. Many of the hands that open $1\cz$ with this shape will have a singleton honour. A typical example is $\sz$K $\hz$KJxx $\dz$AJx $\cz$Q9xxx.

Another possibility with these hands is to pass. This could be a better option than $1\cz$ for hands with 12 HCP or even a bad 13 HCP. Our $1\cz$ opening is already defined to be sounder than $1\dz$ or 1M, and there is no need to open these borderline hands if we believe that any opening bid could lead to problems. If we pass, then there will often be an easy way into the auction later by making a take-out double if the opponents bid our short suit.

\end{description}

\subsection{The `natural' club hands}

The natural option is 15+ HCP with 5+ clubs and no other 5-card suit. Hands with 5 clubs and 4 diamonds should be opened $1\dz$ instead if at the minimum end of this range, but are opened $1\cz$ if strong enough for game opposite a positive response (i.e.~a good 17 HCP or more).

While 15 HCP is usually given as the minimum strength, this is only approximate. Hands with very good club suits can be opened $1\cz$ with less high-card strength. A solid club suit is particularly valuable: a hand containing AKQxxxx in clubs and an outside ace is much too strong for a $2\cz$ bid and must be opened $1\cz$. Hands with 6-4 shape are also frequently upgraded.

When we open $1\cz$ on a `natural' hand with 15-17 HCP, we much prefer to have a good suit. In a competitive auction, often the only way to distinguish the natural type from a weak NT is to rebid the clubs. If our club suit is relatively weak -- and particularly if it is only 5 cards -- we may not be prepared to rebid it at the 3-level. For this reason, we will often consider opening 1NT rather than $1\cz$. Holding 5-4-2-2 shape and a hand in the range for 1NT, we would choose the 1NT opening more often than not. $\sz$KQTx $\hz$xx $\dz$AQ $\cz$KJxxx would be a normal 1NT opener despite the weak doubleton and the lack of rebid problems in an uncontested auction. With 6-3-2-2 shape, a 1NT bid would not be so common, but is still very possible if the hand looks balanced. More difficult are the 5-4-3-1 hands: we should still consider opening 1NT, but the singleton obviously makes this less attractive. $\sz$K $\hz$AQJx $\dz$KQx $\cz$Jxxxx is an easy 1NT opener, but not all hands are this clear. Still it may be best to bid 1NT even on a small singleton in preference to $1\cz$ with a mediocre club suit.

Bidding 1NT on these off-shape hands works best when we are at the lower end of the range for 1NT, because then there is less chance that our extra shape will cause us to miss a game. For hands at the top end of the range it is better to open $1\cz$: if necessary we can pretend to have a `strong' hand, even if our hand does not really evaluate to a full 18 HCP.

With a hand in the `strong' range we don't need to worry about having a good club suit, because there is no need to bid clubs at all.

\subsection{The `strong' hands}

The strong option of $1\cz$ can be any shape: all game-forcing hands are opened $1\cz$. But the shape does matter if we have a hand of 18+ HCP which is not worth a game force. The minimum strength required for the strong option depends on the precise hand type.

\begin{description}
\item[Balanced and 4-4-4-1 hands] With these types, the strong option starts at 18+ HCP -- that is, any hand too strong for a 1NT opening. However, with a 5$\hz$-3-3-2 hand we can choose to open $1\hz$ instead, provided that the hand is no more than 19 HCP.\footnote{You could also agree to open $1\sz$ with a 5$\sz$-3-3-2, though this would increase the upper limit for the $1\sz$ opening slightly.}

\item[Unbalanced hands with 5+ spades]\footnote{With 5-5 shape, we look at the higher-ranking suit (the suit that you would open playing natural methods).} The strong option also starts at a nominal 18 HCP for this hand type. Nearly all 18+ HCP hands with 5+ spades will be opened $1\cz$. Hands with a lower point count may be upgraded to a $1\cz$ opening if the spade suit is good enough, for example $\sz$AKJTxxx $\hz$KJx $\dz$Ax $\cz$x.

\item[Unbalanced hands with 5+ hearts] With a single-suited hand, we open $1\cz$. As with spades, the 18 HCP requirement can be stretched with a particularly good suit. But with a two-suiter, or any hand with precisely 5 hearts, we should open $1\hz$ unless the hand is a game force. 

\item[Unbalanced hands with 5+ diamonds] These must be opened $1\dz$ unless game-forcing.

\item[Unbalanced hands with 5+ clubs] Strong hands with primary clubs are all opened $1\cz$ -- after all, the natural option of the $1\cz$ starts at 15 HCP.

\end{description}

Playing this style, $1\dz$ and $1\hz$ are not particularly limited openings, as they would be in most Polish Club variants. (But note that a $1\dz$ will have at least 5 diamonds if in the `strong' range, and $1\hz$ will not be single-suited.) The main reason we avoid opening $1\cz$ with these hands is because of the possibility of interference. (Also by taking out most hands with diamonds we free up opener's diamond rebids, and avoid the difficult problem of how to bid hands with diamonds after a $1\cz$ opening.) It is particularly difficult to show `flexible' hands after interference over $1\cz$: we could make a take-out double, but this would not tell partner about our 5-card suit. Similarly, two-suiters are difficult to bid in competition if we have not shown a suit with our first bid. But when our longest suit is spades, we are much less worried about interference because we can outbid the opponents on any level; also we have the option of doubling intending to correct partner's response to spades. So it makes a lot of sense for $1\sz$ to be limited while $1\dz$ and $1\hz$ are not. Single-suited hands with hearts are an exception because they are very well described by opening $1\cz$ and rebidding in hearts.

\newpage
\section{The $1\cz$ opening in uncontested auctions}

The responses to $1\cz$ look like this:
\begin{center}
\begin{tabular}{|l|p{10cm}|} \hline
$1\dz$ & \st Usually a `negative', but there are four possibilities:
\begin{itemize}
\item 0-6 HCP, any shape.
\item 6-10 HCP with no 4-card major: balanced or 5$\dz$-4$\cz$.
\item Game force with 5+ diamonds, or invitational with 6+ diamonds.
\item Balanced game force with no 4-card major, not wanting to declare NT.
\end{itemize} \\ \hline
$1\hz$ & \st 6+ HCP with 4+ hearts, may have a longer minor. \\ \hline
$1\sz$ & \st 6+ HCP with 4+ spades, may have a longer minor. \\ \hline
1NT & \st Good 10 or 11 HCP, balanced. \\ \hline
$2\cz$ & \st 6-10 HCP with 5+ clubs. \\ \hline
$2\dz$ & \st 6-10 HCP with 6+ diamonds. \\ \hline
$2\hz$ & \st Forcing to game with 5 clubs (occasionally 6). \\ \hline
$2\sz$ & \st Forcing to game with 6+ clubs \\ \hline
2NT & \st Forcing to game, balanced with no 4-card major. \\ \hline
$3\cz$ & \st 6+ clubs, invitational. \\ \hline
$3\dz$/3M & \st Natural and pre-emptive. \\ \hline
\end{tabular}
\end{center}

\subsection{The negative response}

We call the $1\dz$ response a `negative' even though there are some stronger possibilities included in the bid as well. The four types of hands that bid $1\dz$ are:

\begin{enumerate}
\item[(i)] 0-6 HCP, any shape. This is the genuine negative, and is bid on any hand not strong enough to force to game if opener has the strong type. Good 6 HCP hands, particularly those with a 5+ suit, would normally be worth a positive response instead.
\item[(ii)] 6-10 HCP with no 4-card major, balanced or $5\dz$-$4\cz$. This takes care of balanced hands not strong enough for the 1NT invitational bid. We would usually prefer to bid $2\cz$ with a 5-card club suit.
\item[(iii)] Game-forcing hands with 5+ diamonds, or invitational with 6+ diamonds. This is the only way to show a strong hand with primary diamonds.
\item[(iv)] Game-forcing balanced hands not wanting to declare NT.
\end{enumerate}

When opener has the weak type, he must rebid either $1\hz$ or $1\sz$ (a 1NT rebid would show a strong hand). So the 1M rebids may occasionally be made with only 3 cards. Opener will bid his longer major, and with 4-4 or 3-3 in the majors he would normally bid $1\hz$.\footnote{The $1\hz$ rebid could even be on a 2-card suit if opener has chosen a non-systemic $1\cz$ opening on 2=2=4=5 shape.}

But the 1M rebid is not limited to the weak type: it can also be used with the natural or strong options of $1\cz$. We would start with 1M on any hand with 5+ clubs and a 4-card major, unless it is worth a game force. Similarly, strong hands with 4-4-4-1 shape, or with 5+ spades unsuitable for a jump to $2\sz$, will start with 1M if they are not game-forcing.

Opener's other options are:
\begin{description}
\item[1NT] 18-20 HCP balanced. Strong hands with 5-3-3-2 shape are always treated as balanced after $1\cz$ : $1\dz$, even if they contain a 5-card major.

After the 1NT rebid we play our usual NT system -- this may be a little inefficient because opener will not have a game-forcing hand with a major suit, but it is still worth playing transfers here because of the right-siding effect. A transfer followed by a new suit rebid should be natural and invitational, rather than forcing.
\item[2$\cz$] Natural, 15-22 HCP or so with 5+ clubs; denies a 4-card major. (So, if opener has only 5 clubs, he will have 4 diamonds as well.)
\item[2$\dz$] Artificial game force (24+ HCP if balanced).
\item[2$\hz$/2$\sz$/3$\cz$] 6+ cards, single-suited. An `Acol Two' type hand, but not absolutely forcing.
\item[2NT] 21-23 HCP balanced. Stayman and transfers apply in whatever way the partnership would usually play over a natural 2NT bid.
\item[3$\dz$/3$\hz$/3$\sz$/4$\cz$] Forcing to game, showing a two-suited hand (5-5 or better). $3\dz$ shows both diamonds and a black suit, $3\hz$ shows hearts and a minor, $3\sz$ shows both majors and $4\cz$ shows spades and clubs.
\end{description}

\subsection{Continuations after $1\cz$ : $1\dz$ , 1M}

The 1M rebid is not forcing, but because opener could still have the strong option of $1\cz$, responder will keep the bidding open even with a genuine negative if he has some values (say 4 HCP or more). With 4 or more cards in partner's major he can raise to the 2-level -- this implies about 4-6 HCP. Bidding $1\sz$ over $1\hz$ also shows a negative hand, and promises 4+ spades. And if responder has a very good major suit, he may bid $2\hz$ over $1\sz$ or $2\sz$ over $1\hz$ (but he cannot bid a minor suit at the 2-level, because this would promise a stronger type of hand). If unable to make a suit bid, but with a hand not so bad that it must pass, responder will bid 1NT. However this bid does not necessarily show negative values, since we also bid 1NT on a balanced hand with up to 10 HCP.

Over $1\cz$ : $1\dz$ , 1M : 1NT  or $1\cz$ : $1\dz$ , $1\hz$ : $1\sz$ opener will normally pass with the weak type, or bid 1NT over $1\sz$ with only a doubleton spade. He can show a stronger hand as follows:
\begin{description}
\item[Hands with 5+ clubs.] To show clubs opener rebids either $2\cz$ or $2\dz$. The $2\cz$ bid is limited to about 18 HCP and is not forcing, whereas $2\dz$ is artificial showing 18+ HCP and is forcing for one round. Both of these bids imply 5+ clubs and 4 cards in the major that was bid on the previous round. After $2\dz$, responder can bid 2M or 2OM with a weak hand (the latter being artificial, including hands that want to get out in $3\cz$); everything else is game forcing.

After $1\cz$ : $1\dz$ , $1\hz$ : $1\sz$ we might choose to pass or bid 1NT instead with the natural type, since partner is limited to 6 HCP. Conversely, it would also be acceptable to bid $2\cz$ on the weak type, since we know partner cannot have enough strength to get carried away. But if responder's rebid was 1NT, opener's $2\cz$ promises the natural option of $1\cz$ (at least 15 HCP), since responder can have up to 10 HCP and may want to try for game.

A $3\cz$ bid is \emph{not} an option with the natural type: this bid would suggest that opener's major is at least as long as clubs. Good hands with 5+ clubs all have to start with $2\dz$ here.
\item[Strong hands with 5+ spades.] After $1\cz$ : $1\dz$ , $1\sz$ : 1NT, opener can rebid $2\hz$ to show 5+ spades and 4 hearts, or $2\sz$ to show any other shape with 5+ spades. These rebids are non-forcing. If opener has a hand worth a game force after responder's 1NT bid, he can rebid a second suit naturally at the 3-level.

Note that if opener has 5 spades and a 4-card minor, he cannot bid 2m here as that would show a hand with 5+ clubs and 4 spades, as described above. The 5$\sz$-4m hands must bid $2\sz$.
\item[Strong 4-4-4-1 hands.] After $1\cz$ : $1\dz$, these hands rebid a 4-card major, bidding $1\hz$ with both majors. Now if responder bids 1NT:
\begin{itemize}
\item Opener's 2NT or 3NT rebid shows a 4-4-4-1 with shortage in the other major.
\item $1\cz$ : $1\dz$ , $1\hz$ : 1NT , $2\hz$ shows 4=4=1=4 shape, not forcing.
\item $1\cz$ : $1\dz$ , $1\hz$ : 1NT , $2\sz$ shows 4=4=4=1 shape, not forcing.
\item $1\cz$ : $1\dz$ , $1\hz$ : 1NT , 3m shows 4-4-4-1 with shortage in the other minor, forcing to game.
\end{itemize}
Notice that these sequences after $1\cz$ : $1\dz$ , $1\hz$ are free to show 4-4-4-1 hands because hands with 5+ hearts would either have opened $1\hz$, or, if single-suited, have opened $1\cz$ and rebid $2\hz$.

After $1\cz$ : $1\dz$ , $1\hz$ : $1\sz$, the jump to 2NT or 3NT to show a singleton spade still applies, but with a 4-4-4-1 and both majors we would probably bid 4m as a splinter. A raise to $2\sz$ here would imply the weak type instead, though in practice you are unlikely to have an uncontested auction if holding those hands.
\end{description}

Responder's $2\cz$ and $2\dz$ rebids after $1\cz$ : $1\dz$ , 1M are both artificial:
\begin{itemize}
\item $2\cz$ shows the strong hand with 5+ diamonds. Opener is asked to bid $2\dz$ if he has the weak type.
\item $2\dz$ shows 6-10 HCP with 5 diamonds and 4+ clubs. Opener will pass or bid 3m with the weak type.
\end{itemize}
After either of these bids, we have an artificial way for opener to show the natural option of $1\cz$: a $2\hz$ rebid shows this type of hand. So this bid is at least invitational after $2\dz$ (the bidding can stop in 3m), and is game-forcing after $2\cz$. In fact we will see this artificial $2\hz$ bid in various other places in the system: whenever responder shows diamonds at the 2-level, opener's $2\hz$ bid shows the natural option of $1\cz$.

If opener has a strong hand after $1\cz$ : $1\dz$ , 1M : $2\cz$/$2\dz$, he could bid a natural $2\sz$, or 3OM as a splinter, or 4m as a natural slam try.

After $1\cz$ : $1\dz$ , 1M : $2\cz$ , $2\dz$ (where responder has shown 5+ diamonds and opener has shown the weak type), responder rebids as follows:
\begin{description}
\item[2$\hz$] Shows precisely 5 diamonds, forcing to game.
\item[2$\sz$] Shows 6+ diamonds, forcing to game.
\item[3$\cz$] Shows 5-5 or better in the minors, invitational.
\item[3$\dz$] Invitational.
\item[3$\hz$/3$\sz$] Forcing to game, 5-5 or better in the minors, shortage in the other major.
\item[3NT] To play -- no interest in other contracts opposite the weak type.
\end{description}
This is not just some arbitrary arrangement: the $2\hz$ and $2\sz$ bids are intended to correspond to the sequences $1\cz$ : $2\hz$ and $1\cz$ : $2\sz$, which show game-forcing hands with \emph{clubs}. The continuations can be found in the section on those bids.

\vspace{4mm}
With the game-forcing balanced hand which does not want to be declarer in no-trumps, responder bids $3\cz$ over 1M. Opener will then bid:
\begin{description}
\item[3$\dz$] Shows shortage in diamonds (any strength). Now $3\hz$ asks about strength (opener will bid $3\sz$ if weak). With opener having a shortage it is likely to be better for responder to declare 3NT if that is to be the final contract, despite not wanting to bid 2NT initially.
\item[3$\hz$] Shows the natural type with major-suit shortage.
\item[3$\sz$/4-any] Show strong hands.
\item[3NT] Any weak hand apart from the $3\dz$ bid.
\end{description}

\subsection{Continuations after $1\cz$ : $1\dz$ , $2\cz$}

Opener will have either 6+ clubs or a hand with 5 clubs and 4 diamonds. When he has 6+ clubs the range will be about 15-20 HCP, whereas with the 4-5 hand it is more like 17-22 HCP.

The continuations are based around a $2\dz$ asking bid. Opener replies to show his strength and club length:
\begin{description}
\item[2$\hz$] Minimum with 6+ clubs (up to about 18 HCP).
\item[2$\sz$] Minimum with 5 clubs and 4 diamonds.
\item[3$\cz$] Maximum with 6+ clubs.
\item[3$\dz$] Maximum with 5 clubs and 4 diamonds.
\end{description}
Opener may also bid 2NT or 3M (showing stoppers) with a super-maximum hand with 5 clubs and 4 diamonds -- just below a game force. With a single-suiter, a hand of this strength would have rebid $3\cz$ over the negative.

Because opener's $2\cz$ bid starts at 15 HCP for single-suited hands but nearer 17 HCP for hands with 5 clubs and 4 diamonds, opener's minimum range is wider when he has 6+ clubs. Over the $2\hz$ rebid, responder can make a further invitation by bidding $2\sz$ (opener will bid $3\cz$ with an absolute minimum).

After any reply to $2\dz$, bids of new suits at the 3-level show stoppers.

The $2\dz$ enquiry can be made on any strength of hand, including some negative hands, nearly all hands of 6-10 HCP, and also with most game-forcing hands (there is no other way to show the game-forcing balanced type). Since $2\dz$ does not promise anything in the way of values, when opener has shown a minimum he must respect partner's 3m sign-off. If responder wants to make a slam try in clubs or diamonds he can bid 4m after finding out about opener's overall strength.

Responder's other bids over $2\cz$ are natural. 2M shows a negative hand with 5+ cards in the suit, contructive but not forcing. $3\cz$ also shows a negative, to play unless partner has an absolute maximum. Jumps to $3\dz$ or 3M show the strong type with 5+ diamonds, and will tend not to have a fit for clubs otherwise we would start with $2\dz$ instead.

\subsection{Opener's game-forcing hands after $1\cz$ : $1\dz$}

Opener's $2\dz$ rebid is absolutely forcing to game. Responder's rebids primarily describe his major-suit length:

\begin{description}
\item[2$\hz$] Waiting, denies holding a 5-card major.
\item[2$\sz$] Shows an unspecified 5-card major, 0-6 HCP.
\item[2NT] Natural, positive values.
\item[3$\cz$] 5 diamonds and 4+ clubs, positive values.
\item[3$\dz$] Natural, 10+ HCP, good suit.
\end{description}

Subsequent bidding is mostly natural. After $1\cz$ : $1\dz$ , $2\dz$ : $2\sz$ , 2NT, transfers apply as normal but $3\cz$ `Stayman' implies 5-4 in the majors. After $1\cz$ : $1\dz$ , $2\dz$ : $2\hz$ , 2NT, normal Stayman is on.

The advantage of responder immediately describing whether he has a 5-card major is that after $1\cz$ : $1\dz$ , $2\dz$ : $2\hz$ , $3\cz$/$3\dz$ responder can bid a major when nothing else appeals without partner expecting him to have a real suit.

\vspace{3mm}
Opener also has special rebids available over $1\cz$ : $1\dz$ which show game-forcing two-suited hands, 5-5 or better. Specifically:
\begin{description}
\item[3$\dz$] Diamonds and a black suit. Responder bids next step ($3\hz$) to ask for the second suit.
\item[3$\hz$] Hearts and a minor. Responder bids 3NT to ask for the minor.
\item[3$\sz$] Both majors. Responder can now bid $4\cz$ to make a slam try in hearts, $4\dz$ for a slam try in spades.
\item[4$\cz$] Spades and clubs.
\end{description}
So in principle opener's artificial $2\dz$ game force will not be made on 5-5 shape, but with a bad second suit it is a possibility, particularly holding spades and a minor, where we will normally be able to bid spades at the 2-level.

Note that these artificial jumps do not apply in competition, unless all the opponents have done is double.

\subsection{The major-suit responses}

Opener's $1\hz$ or $1\sz$ response shows 4+ cards and promises enough strength for game if opener has the strong type -- so at least a good 6 HCP. As in most systems, we bid $1\hz$ with 4-4 in the majors, but $1\sz$ with 5-5.

Responder could have a longer minor suit. With less than a game force, hands with 5+ diamonds and a 4-card major must bid 1M, whereas hands with 5+ clubs and a 4-card major have a choice between 1M and $2\cz$. On a game-forcing hand, we would normally show the minor first unless we have 5-4-2-2 shape, in which case we must bid the major -- the system does not give us a way to show this semi-balanced hand if we start with the minor.

Opener rebids as follows:
\begin{description}
\item[With the weak type.] Opener will raise 1M to 2M with 4-card support, or with an unbalanced hand with 3-card support (4=3=1=5 or 3=4=1=5). Otherwise he will rebid 1NT, or possibly $1\sz$ over $1\hz$. Recommended is to bid 1NT over $1\hz$ on any balanced hand, even with 4 spades; then $1\sz$ promises 5+ clubs and 4 spades.\footnote{As is the case in natural systems, there is an alternative (more traditional) style of bidding $1\cz$ : $1\hz$ , $1\sz$ on some or all balanced hands with 4 spades. This is perfectly compatible with Polish Club, and some players might prefer to play it this way. This choice is linked to the checkback methods used over 1NT.} After either of these rebids we are essentially in the same position as we would be if we had opened $1\cz$ in a natural system. So after the 1NT rebid we should play some sort of checkback (a `transfer checkback' scheme is explained in the section on optional methods), whereas after $1\cz$ : $1\hz$ , $1\sz$ we can play natural-based continuations with $2\dz$ as fourth-suit forcing.

Note that if we opened $1\cz$ on a weak type with 5 clubs and 4 hearts, but without 3 spades, we must always rebid 1NT over a $1\sz$ response. So the 1NT rebid could occasionally be made on a singleton with precisely 1=4=3=5 shape. We cannot rebid $2\cz$ because this would show a stronger hand.

\item[With the natural type.] Over $1\cz$ : $1\hz$, hands with 5+ clubs and 4 spades can rebid $1\sz$. This rebid should be forcing: responder will have at least a good 6 HCP, and so is unlikely to want to pass.

Most other natural hands will rebid $2\cz$, which is forcing for one round (see continuations below). If the hand is worth a game force after respodner's 1M bid then there are other possible options: a game-forcing hand with 3-card support for partner should rebid $2\dz$ instead, and a game-forcing single-suited hand can rebid $3\cz$ (implying an `Acol Two' type). But most of the weaker hands, along with many game-forcing ones, must start with a $2\cz$ bid.

With 4-card support we can raise to 3M (invitational), or bid 3OM or $4\dz$ as a splinter. A jump to $4\cz$ shows 4-card support with 6+ good clubs.

\item[With the strong type.] Opener's $2\dz$ rebid is an artificial game force: it shows either 18-20 HCP balanced, or any strong hand with 3-card support for partner's major (but hands with a very good suit of their own can bid that suit instead).

Opener's 2NT rebid shows 21+ HCP and denies 3-card support. Continuations are natural. Responder's minor-suit rebids could be 5-4 shape with either suit being longer; after this opener's 3M rebid shows a suitable hand for playing in a suit contract, with support for the minor. Similarly in the sequence $1\cz$ : $1\hz$ , 2NT : $3\sz$ we do not know whether responder's hearts are longer than spades: it could be 4-4 or 4-5.

Hands with 4-4-4-1 shape short in partner's suit have to be treated as balanced.

Bidding the other major at the 2-level ($1\cz$ : $1\hz$ , $2\sz$ or $1\cz$ : $1\sz$ , $2\hz$) shows a strong hand with 5+ cards in that major.\footnote{An alternative is to allow opener to rebid $1\sz$ over $1\hz$ on strong hands with spades; however this makes the $1\cz$ : $1\hz$ , $1\sz$ sequence rather complicated.} With the very strong hand with diamonds we have to jump to $3\dz$ (unless we have 3-card support). Again the bidding should continue naturally.

Note that a jump to 3 of the other major is not natural -- this is a splinter with 4-card support for partner, usually with length in clubs. However a jump to game is natural, suggesting a minimum strong hand with a particularly long suit. And to make a splinter in clubs (for example with 4=4=4=1 shape) we have to start with $2\dz$.
\end{description}

\subsection{$1\cz$ : 1M , $2\cz$}

This is forcing, however opener has other options on a game-forcing hand with clubs: with a single-suiter he can bid $3\cz$, and with support for partner he can bid $2\dz$ or make a splinter.

Responder rebids as follows:

\begin{description}
\item[2$\dz$] Any weak hand except for a 2M bid. Now opener's rebids of 2M or 2NT or $3\cz$ are non-forcing, as is the sequence $1\cz$ : $1\sz$ , $2\cz$ : $2\dz$ , $2\hz$. Opener's rebids of new suits above 2M are forcing.

Responder would also bid $2\dz$ on a minimum game force with only 4 cards in his major, if he is weak in one of the unbid suits. (With honours in both unbid suits he can bid 2NT instead.) With this hand he will intend to rebid a new suit at the 3-level. Other rebids show weaker hands.
\item[2M] Shows a good suit, 6+ cards in principle, constructive but not forcing.
\item[2OM] Artificial game force promising 5+ cards in the original major.
\item[2NT and higher] Natural and forcing to game. A new suit at the 3-level promises at least 5 cards there ($3\dz$ is likely to be 4M-$5\dz$ shape; 3OM will be at least 5-5). A $3\cz$ bid encourages opener to make a slam try if he has extras -- if responder has a game force which is a bad hand for slam he could start with $2\dz$ or 2NT instead.
\end{description}

\subsection{$1\cz$ : 1M , $2\dz$}

As in most versions of Polish Club, the $2\dz$ rebid is an artificial game force. We use it to show either 18-20 HCP balanced or any strong hand with 3-card support for partner's major. Responder rebids:

\begin{description}
\item[2$\hz$] Precisely 4 cards in the major. Now opener bids 2NT with 18-20 HCP balanced, and $2\sz$ with 21+ HCP balanced (implying 3-card support). 3-level bids are natural and also promise 3-card support.
\item[2$\sz$] Precisely 5 cards in the major. Again opener bids 2NT with 18-20 HCP balanced, even if holding 3-card support for responder's suit. Stronger balanced hands bid 3M setting trumps.
\item[2NT] 6+ cards in the major, minimum.
\item[3$\cz$/3$\dz$/3OM] Natural, at least 5-5.
\item[3M] 6+ cards in the major, better than minimum.
\end{description}

Notice that when responder makes the most common $2\hz$ and $2\sz$ rebids, he does not immediately describe the strength of the hand. But opener will be able to show his strength instead (at least when he has a balanced hand, which is the most common situation), after which responder can decide whether to make a slam try.

After $1\cz$ : 1M , $2\dz$ : $2\hz$ , $2\sz$/2NT, responder's bid of a new suit is natural with 3m promising 5+ cards. Rebidding his major is artificial: when opener bid $2\sz$, its main function is asking opener to bid 3NT in order to right-side the contract; however it is also the way to start a slam try on a hand without a 5-card minor.

\subsection{The 1NT and 2NT responses}

Responses in no-trumps are natural. 1NT is invitational opposite the weak type, and so promises a good 10 or 11 HCP, whereas 2NT is a game force on 12+.

The 1NT bid does not deny a 4-card major, but we would prefer to respond 1M if finding a major-suit fit would make a significant difference to the evaluation of the hand. So 1NT will tend to have a 4-card major only when 4-3-3-3, or with a 4-4-3-2 which looks particularly suited to no-trump contract: such hands do not increase much in value when a fit is found. The 2NT response, on the other hand, should in principle deny a 4-card major.

Over 1NT, any rebid from opener is forcing to game (hands wanting to decline the invite will pass). The rebids are:
\begin{description}
\item[2$\cz$] Shows the natural option of $1\cz$.
\item[2$\dz$] Shows the weak option with shortage in diamonds.
\item[2M/3m] Natural strong hands.
\item[2NT] Balanced (or a weak type with 5 clubs), but interested in a suit contract. This is a puppet to $3\cz$: responder must bid $3\cz$, and then opener normally bids a 4-card suit at the 3-level. Responder will raise holding 4-card support, and bid 3NT otherwise. Alternatively, opener can rebid 3NT over $3\cz$ to show both majors and no slam interest, or $4\cz$ to show a balanced slam try with clubs.
\item[3NT] To play.
\end{description}

Over 2NT, things work in a similar way, but there is no `Stayman' bid here.
\begin{description}
\item[3$\cz$] Shows the natural option of $1\cz$.
\item[3$\dz$] Shows the weak option with shortage in diamonds.
\item[3M/4m] Natural strong hands.
\item[3NT] To play.
\end{description}

\subsection{The $2\cz$ response}

This is natural, and essentially `to play' opposite a weak NT, but opener is allowed to make a natural raise with a good fit for clubs even if minimum. $2\cz$ does not deny a 4-card major, but particularly at the maximum end of the range we would normally prefer a 1M bid.

Apart from natural raises, opener's rebids are all forcing to game. It is possible that we might not have the high-card strength for game if opener has the natural type, but with an unbalanced hand and a 10-card fit this seems unimportant, and in any case it would be very rare to hold this hand with the opponents both passing.

Opener's $2\dz$ rebid is artificial, showing good club support. This includes the natural option of $1\cz$, but could also be any strong hand with 3+ clubs.

Strong hands with a 5+ major can rebid 2M. With a game force in diamonds we have to jump to $3\dz$, which is natural; however jumps to 3M are splinters.

A 3NT bid is to play, whereas 2NT is likely to be a stronger hand or considering a $5\cz$ contract. After opener rebids 2NT, responder's rebids are conventional:
\begin{description}
\item[3$\cz$] Waiting, starts a natural sequence.
\item[3$\dz$] Shows shortage in diamonds.
\item[3$\hz$] Shows shortage in spades.
\item[3$\sz$] Shows shortage in hearts.
\item[3NT] No shortage, minimum hand.
\end{description}
The same scheme will be used throughout the system in situations where responder has shown a 5-card minor and opener bids 2NT showing the strong option. (When responder has shown diamonds, we still use $3\cz$ as the waiting bid, and $3\dz$ will show shortage in clubs.) In these situations 3M always shows shortage in the opposite major.

\subsection{The $2\dz$ response}

This is similar to the $2\cz$ response in that responder expects partner to pass or raise with the weak NT type, but here we need a better suit because opener could have shortage in diamonds.

After $2\dz$, opener's $2\hz$ rebid is artificial and shows the natural option of $1\cz$ (including hands in the strong range). This bid is therefore at least invitational opposite responder's 6-10 HCP range. The bidding can stop in $3\cz$ or $3\dz$, but otherwise is forced to game.

So when responder has a strong hand with hearts, he must jump to $3\hz$. This is not a splinter like it would be over a $2\cz$ response. However a jump to $3\sz$ is still a splinter, since we can bid $2\sz$ with natural spades.

Opener's other rebids are all natural.  He can bid $3\dz$ with the weak type; otherwise he is showing a strong hand. $3\cz$ would show a game-forcing single-suiter. Like over the $2\cz$ response, a jump to 3NT is to play, whereas 2NT is natural and forcing with conventional replies. Following the same pattern as over the $2\cz$ response, responder's rebids after 2NT are:
\begin{description}
\item[3$\cz$] Waiting, starts a natural sequence.
\item[3$\dz$] Shows shortage in clubs.
\item[3$\hz$] Shows shortage in spades.
\item[3$\sz$] Shows shortage in hearts.
\item[3NT] No shortage, minimum hand.
\end{description}

\subsection{Responses showing good hands with clubs}

$2\hz$ and $2\sz$ are artificial game-forcing bids showing clubs; usually we would bid $2\hz$ with precisely 5 clubs and $2\sz$ with 6+, but 6-3-2-2 hands with poor clubs could go through $2\hz$ instead. Note that hands with 5-4-2-2 shape and a 4-card major must respond 1M rather than $2\hz$: the continuations after $2\hz$ do not give us a way to show this hand.

After either $2\hz$ or $2\sz$, opener will usually rebid 2NT with the weak type, but can also bid $3\cz$ with a particularly good club fit. The 2NT rebid also takes care of most strong balanced hands, except that a minimum strong type with poor support for clubs would jump to 3NT. Other rebids are natural strong hands.

Responder's rebids are natural except when opener bid 2NT. Here the rebids follow the same general pattern as over $1\cz$ : 2m , 2NT, but the hand types are more specific. Again the emphasis is on showing shortage.

\begin{center}
\begin{tabular}{|l|p{4.5cm}|p{5cm}|} \hline
 & \st Over $1\cz$ : $2\hz$ , 2NT & Over $1\cz$ : $2\sz$ , 2NT \\ \hline
$3\cz$ & \st Shows 4 $\dz$s & Shows a 6-4 shape with an unspecified 4-card suit \\ \hline
$3\dz$ & \st Diamond shortage & Diamond shortage \\ \hline
$3\hz$ & \st Spade shortage, with 4 $\hz$s & Spade shortage \\ \hline
$3\sz$ & \st Heart shortage, with 4 $\sz$s & Heart shortage \\ \hline
3NT & \st To play opposite the weak type & To play opposite the weak type \\ \hline
$4\cz$ & 5-3-3-2 slam try \st & Slam try with no shortage \\ \hline
\end{tabular}
\end{center}
\vspace{3mm}
After $1\cz$ : $2\sz$ , 2NT : $3\cz$ opener rebids $3\dz$ to ask for the 4-card suit, with 3M showing 4 cards there and 3NT/$4\cz$ showing 4 diamonds. Similarly after $1\cz$ : $2\hz$ , 2NT : $3\cz$ opener can ask for clarification with $3\dz$: then $3\hz$ shows 1=3=4=5, $3\sz$ shows 3=1=4=5 and 3NT shows 2=2=4=5.

\vspace{4mm}
Exactly the same scheme is used in the sequence $1\cz$ : $1\dz$ , 1M : $2\cz$ , $2\dz$ : $2\hz$/$2\sz$ , 2NT, where responder has shown a game forcing hand with \emph{diamonds}. Bids which showed features in clubs now show diamonds, and vice versa, so for example
\begin{itemize}
\item $1\cz$ : $1\dz$ , 1M : $2\cz$ , $2\dz$ : $2\hz$ , 2NT : $3\cz$ shows 5 diamonds and 4 clubs.
\item $1\cz$ : $1\dz$ , 1M : $2\cz$ , $2\dz$ : $2\hz$/$2\sz$ , 2NT : $3\dz$ shows shortage in clubs.
\end{itemize}

\vspace{3mm}
The $3\cz$ response to $1\cz$ is natural and invitational. After this, opener's $3\dz$ or $3\hz$ rebids initially show stoppers (but could be the first move on a stronger hand), whereas a $3\sz$ rebid will always be 18+ with 5+ spades. To make a slam try in clubs opener can raise to $4\cz$; opener's jump to 4NT is natural.

\subsection{When responder is a passed hand}

We do not change any of the system for uncontested auctions when $1\cz$ was opened in third or fourth seat, though opener might choose to pass a 1M response.

The game-forcing responses might seem unlikely, but there is not really any useful alternative meaning for those bids. With both opponents having passed we would only be pre-empting our partner, who is quite likely to have a strong hand in this situation. And it is just about possible for responder to have a $2\hz$ bid, since we occasionally pass difficult 12- or 13-point hands with 5 clubs.

\newpage
\section{The $1\cz$ opening in competition}

Assume for the moment that the opponents' interference is natural. Many opponents will be playing artificial methods, and we will come to dealing with artificial interference later, but first we need to know what happens over natural overcalls and doubles.

Let's start with some general principles for responder's bidding after opener's LHO has made a natural overcall:

\begin{itemize}
\item Any action apart from a pass creates a game force if opener has the strong type.
\item Double is for take-out.
\item 1-level bids are natural, with 1-of-a-suit being forcing by an unpassed hand.
\item 2-level suit bids are generally natural and non-forcing.
\item BUT after a $1\sz$ overcall we play special methods, with transfers starting at $2\cz$.
\item After 2-level interference, we play transfers starting at 2NT.
\item The cue-bid of opponents' suit is given some artificial meaning (e.g.~a 2-level cue-bid shows a good hand with clubs).
\end{itemize}
More on all of this later.

\subsection{Opener's reopening actions}

Suppose that opener's LHO makes a natural overcall, and responder passes. The pass will either be a hand too weak for any action, or a decent hand with some length in the opponents' suit. Now RHO also passes, and it is back to opener.

With the weak option of $1\cz$, opener usually has to pass it out, but at low levels he has the option of making a take-out double:

\begin{itemize}
\item After a suit overcall up to and including $2\dz$, opener may make a reopening double with a suitable minimum.
\item After an overcall of $2\hz$ or higher, opener's reopening double promises at least the natural option of $1\cz$: opener must pass with the weak type.
\end{itemize}

With the natural option, opener can double for take-out or bid clubs naturally. If he has clubs and a major he must not bid the major: this would show the strong type and promise at least 5 cards. (There is one exception: reopening with a bid of $1\hz$ -- which can only happen after a $1\dz$ overcall -- shows clubs and hearts, because hands with 5 hearts would open $1\hz$ and hands with 6+ hearts can afford to bid $2\hz$.) Of course, opener is not obliged to bid again with the natural type and may choose to pass, particularly with length in the opponents' suit and at high levels.

Other bids show the strong type (though again, doubling for take-out is an important option on strong hands). Bidding is generally very similar to if the opponent's bid was an opening bid rather than an overcall. A 1NT or 2NT bid is always natural, promising the strong type even at the 1-level, with a jump being stronger than a simple overcall. A jump to 3NT, however, is more likely to be a hand with long clubs.

With a flexible hand with 5 spades, we would normally double intending to bid spades on the next round, rather than bidding spades immediately (except at the 1-level, where $1\sz$ shows this hand nicely). So a suit bid tends to imply a single-suited hand. Also note that a natural bid in diamonds is forcing, because hands not worth a game force would have opened $1\dz$ instead; similarly a double followed by a bid in hearts is logically forcing.

There are some important artificial rebids for opener:

\begin{itemize}
\item A rebid of $2\dz$ is an artificial game force, exactly the same as $1\cz$ : $1\dz$ : $2\dz$, and with the same continuations. (So $2\hz$ now from responder denies a 5-card major and $2\sz$ shows one.) 
\item A cue-bid of the opponents' suit below 3NT shows clubs. This is essentially the same as the common defence to weak twos whereby ($2\hz$) : $3\hz$ shows a strong single-suited hand, except that for us it shows specifically clubs since hands with diamonds can bid a forcing $3\dz$. When the cue-bid is a jump it would imply a solid suit.
\item When the opponent's overcall was 1M or 2M, a jump to 4 of a minor is `Leaping Michaels' showing a two-suited hand with the other major, at least 5-5, and forcing. Again, this is a normal method for dealing with weak twos, and it makes a lot of sense for us to use it over overcalls of our $1\cz$ as well. Similarly a jump to 4M shows both minors.
\end{itemize}

Another convention for dealing with weak two openings that we apply here is Lebensohl. After  $1\cz$ : (2x) : Pass : (Pass) , Dbl, a 3-level bid from responder shows values and creates a game force if opener has the strong type. With a weaker hand responder has to go through 2NT -- opener will normally rebid $3\cz$, and then responder will bid a suit to play. Because opener's strong hands start at 18+, the requirement for bidding at the 3-level is slightly lower than over a double of a weak two opening: any hand which is happy to go to game when opener has the strong option is good enough.\footnote{The Lebensohl 2NT can also include some hands which are worth a game force opposite opener's double: I like to play that 2NT followed by 3NT shows doubt about 3NT, whereas 2NT followed by a cue is `Stayman' with a stop in opponents' suit. But I'm often told that this is inconsistent. Whatever you like to do, it makes sense to play these things exactly the same way that you would over a double of a weak two opening.}

Occasionally opener may have to pass with the strong type, but this is very rare: it would only happen at high levels (probably above 3NT) when we have length in opponents' suit. At high levels, double is essentially `value-showing' and most balanced hands are suitable; still opener should be prepared for partner to take out the double on a particularly distributional hand, and so it is better to pass if that would be disastrous.

\subsection{Interference from fourth seat when responder has not shown values}

This applies after a negative response, or if opener's LHO made an overcall and responder passed. In these situations, methods over fourth-seat interference are very similar to those for reopening actions. So $2\dz$ is again an artificial game force, cue-bids show clubs, and generally methods are similar to if RHO's bid was an opening bid rather than an overcall. But an important difference when we are not in the protective position is that opener must pass with the weak type: a double promises a stronger hand.

There is just one difference in system compared to the previous section. When opener doubles RHO's 2M bid, ordinary Lebensohl no longer applies. This is because responder is unlikely to have a hand good enough for game opposite the strong type. Instead we use a form of `scrambling' 2NT. This 2NT bid is used on three different types of hands:
\begin{itemize}
\item Weak hands unsure of which suit to play in (usually `two places to play'). Partner will choose a minor; if we have hearts and the other minor (over a double of $2\sz$) then we will correct.
\item Very weak hands wanting to take out into the suit below the one doubled.
\item The same game-forcing hands that would go through 2NT playing Lebensohl. We will cue or bid 3NT later (or make some other bid if these are defined) as appropriate.
\end{itemize}
Hands which do not need to scramble can just bid their suit. A bid at the 3-level here does \emph{not} promise enough strength for game opposite the strong type, and so opener will pass with a minimum strong hand. However, bidding the suit below the one doubled is slighly more constructive, though still not game-forcing opposite the strong type, showing perhaps 4-5 HCP: with weaker hands we would go through 2NT (the second of the options above), and so 2NT does act as `Lebensohl' on those hands.

We will see leter that a similar 2NT bid is used by \emph{opener} in replying to responder's take-out doubles.\footnote{This `scrambling Lebensohl' and is a useful improvement on the standard scrambling 2NT, and not only after a Polish $1\cz$ opening. Particularly after a take-out double of $2\sz$, it is useful to have a constructive $3\hz$ bid available.}

\vspace{4mm}
If a $1\dz$ negative response is doubled, we play mostly `system on'. However, opener can pass with a weak hand and 4+ diamonds if he does not wish to bid 1M. Also, redouble is an option on a strong hand.

\subsection{Interference from fourth seat after a positive response}

For the moment we can assume that opener's LHO passed, but in fact the methods in this section will almost always apply when responder was bidding over an overcall or a double as well.

When responder has shown a major suit, opener's double of a suit overcall at the 1- or 2-level is a \emph{support double}, promising 3-card support for partner's suit. The support double is not compulsory with 3-card support, but we would normally only hide the support with a very defensive minimum hand, or with a strong hand which has a more descriptive bid available (for example with a long suit which opener wants to show instead). Support doubles apply even if the opponents' bid is artificial.

A sequence such as $1\cz$ : (Pass) : $1\hz$ : ($2\sz$) , Dbl is a special case because responder may be forced to bid but cannot return to his suit at the 2-level. In this one situation, which only occurs when a heart bid is overcalled with $2\sz$, the support double promises at least the natural option of $1\cz$ (opener must pass with the weak type or bid $3\hz$ with a very suitable hand). Here responder will bid $3\cz$ if he wants to play there opposite the natural type, and so this bid does not promise a genuine club suit.

All other doubles of suit bids are for take-out, though tending to become more `value-showing' at higher levels. A double at the 3-level will typically be a strong hand without a stop.

Apart from the support double, opener's options with the weak type are:
\begin{itemize}
\item Pass. This is the normal action with a weak hand.
\item Make a direct raise, implying 4-card support. However, if partner's bid was only at the 1-level, it would be rare to raise to the 3-level with the weak type, and opener should never raise to the 4-level (this ought to promise a stronger hand).
\item Bid 1NT if available. Note that it is only after a positive response that opener can bid a `free' 1NT on a weak hand. But it may be better to pass instead, depending on the hand. 1NT should show a decent opener, and a good stop in the opponents' suit if they have bid one.
\end{itemize}

Similarly, if a 1M response is doubled, opener's redouble shows 3-card support. Minimum hands without 3-card support can pass or bid 1NT (or $1\sz$ over $1\hz$) depending on whether we want to be declarer in NT.

\vspace{4mm}
As for opener's stronger hand types, we try to play `system on' as far as possible, but much of the time the opponents' intervention will have taken our system away. But we have methods to deal with this, including frequent use of diamond bids as artificial. The following applies after any positive response, not just 1M.

\begin{description}
\item[`System on' for $2\dz$/$2\hz$ rebids.] If a $2\dz$ rebid is available, it shows the same hand is it would without the intervention. So if partner bid a major, $2\dz$ shows either 18-20 HCP balanced or 3-card support. If partner bid an invitational 1NT, $2\dz$ shows shortage in diamonds. (If partner bids 1NT showing 7-11 HCP in competition, then $2\dz$ is game-forcing with clubs). And if partner bid clubs, $2\dz$ shows club support.

Similarly, when responder shows diamonds at the 2-level, opener's $2\hz$ bid is artificial and shows the natural option of $1\cz$.\footnote{When opener's LHO passed, this can only happen if a $2\dz$ response is doubled, but we shall see later that there are more general situations where this applies.}

\item[With the natural type.] If $2\cz$ is available, this is natural and shows the same hand as it would without the intervention. So if partner showed a major then this is forcing, with the usual system over $1\cz$ : 1M , $2\cz$ applying. A jump to $3\cz$ would be a game-forcing single-suiter, like in an uncontested sequence.

More often, RHO's intervention takes away our $2\cz$ bid, so we need to use a completely different system. When we have no way to show clubs below $3\cz$, a non-jump $3\cz$ bid is natural and \emph{not} forcing; it also denies 3-card support for partner because we would be able to double with that. To make a forcing bid with clubs we may have to start with a double (either as a support double or as an artificial force at the 3-level) or a cue-bid. However, more often we will be able to bid $3\dz$ (see below).

\item[The $3\dz$ convention.] Unless partner has shown diamonds, a non-jump $3\dz$ rebid is artificial and forces to game. (Indeed, this applies to \emph{any} non-jump $3\dz$ rebid in competition when partner has shown values, except when he made a take-out double.) The precise meaning of $3\dz$ depends on whether some other cue-bid is available below 3NT. If there is no other cue-bid available then the $3\dz$ bid acts as a general forcing bid, as a substitute for the cue. However, if there is another cue-bid available (say opponents bid spades naturally, so that our $3\sz$ would be a cue), then the $3\dz$ bid shows specifically clubs. So generally this shows the natural hand type (but too strong for a non-forcing $3\cz$), but if partner's suit is clubs, it just shows a strong hand with good support.

\item[With a strong balanced hand.] 2NT and 3NT are natural strong bids, and these would normally deal with most strong balanced hands with some sort of stopper. 3NT is to play (unless partner has significant extras) whereas 2NT is more interested in other contracts, including a possible slam try. If partner bid a major, we would normally prefer to make a support double if possible, so no-trump bids deny 3-card support in principle. The 2NT bid can be made with only half a stop in the opponents' suit: responder will be able to check how good the stop is.

If responder's suit was a major, then continuations after 2NT are natural, the same as after $1\cz$ : 1M , 2NT (apart from the cue, which can be used to check whether we have a real stopper). However, if responder's suit was a minor, we use the same artificial continuations as over $1\cz$ : 2m , 2NT. That is, $3\cz$ is waiting and starts natural bidding (including hands which need to check on stoppers), $3\dz$ shows shortage in the other minor, and $3\hz$ or $3\sz$ shows shortage in the opposite major.

Also if partner bid a major, we may have a $2\dz$ bid available. This takes care of the 18-20 HCP hands (so that 2NT promises 21+) as well as being a better option than a support double on hands with 3-card support. However, if we bid $2\dz$ on a balanced hand then we should have a stop in the opponents' suit, since we will have to rebid in no-trumps.

Lacking a stop in the opponents' suit, we have to double or make a cue-bid.

\item[With a strong unbalanced hand.] We show a strong hand with 5+ cards in an unbid major by bidding it at the lowest available level (except that when partner has shown diamonds, $2\hz$ is an artificial bid, and so we may need to jump to $3\hz$ to show hearts).

If our long suit is diamonds, we may not be able to bid the suit naturally immediately because diamond bids are conventional. So in order to show diamonds we have to first make an artificial game-forcing bid (or double) and then rebid in diamonds.

\item[With 4-card support.] After $1\cz$ : 1M in an uncontested auction, strong balanced hands with 4-card support would rebid $2\dz$. But the interference may have taken this option away. Now we might make a natural raise to 4M, which promises a strong hand. However, sometimes we may have a hand too strong for a non-forcing 4M. In this case we use an artificial $4\dz$ bid to show a serious slam try in partner's suit (better than a minimum strong type).

When partner bid a minor, a natural 4m raise shows a serious slam try.

We may also have splinter bids available: a jump to 3 of an unbid major is a splinter, except for $3\hz$ in the situations where $2\hz$ would be artificial. And unless the opponents' suit is diamonds, a 4-level cue-bid is also a splinter. So for example after $1\cz$ : (Pass) : $1\sz$ : ($3\hz$), opener's $4\dz$ rebid is a general slam try in spades while $4\hz$ shows specifically shortage in hearts.
\end{description}

\subsection{Responder's 2-level non-forcing suit bids}

After interference, responder's 2-level suit bids are mostly played as natural and non-forcing. For example:

\begin{itemize}
\item $1\cz$ : ($1\hz$) : $2\cz$
\item $1\cz$ : ($1\dz$) : $2\sz$
\item $1\cz$ : (1NT) : $2\dz$
\item $1\cz$ : (Dbl) : $2\hz$
\item $1\cz$ : ($2\cz$) : $2\hz$
\item $1\cz$ : ($2\hz$) : $2\sz$
\end{itemize}
We will refer to these as `competitive free bids' (CFBs). It makes no difference whether the bid was a jump, or what level the opponents' interference was at. The important exception is after a $1\sz$ overcall: here natural CFBs do not apply (we play transfers instead). Apart from this, responder's 2-level suit bid in competition is always a CFB unless it is a cue-bid of the opponents' suit.

Note that there are two uncontested sequences which really belong here as well:

\begin{itemize}
\item $1\cz$ : (Pass) : $2\cz$
\item $1\cz$ : (Pass) : $2\dz$
\end{itemize}
Our CFBs in competition will be treated in the same way as these minor-suit bids. (But in a competitive auction we will often not have such as good a suit as is required for the uncontested $1\cz$ : $2\dz$.)

\vspace{4mm}
These non-forcing bids are an essential part of our bidding after a $1\cz$ opening. They are similar to `negative free bids' after a natural suit opening, perhaps not promising quite as good a suit. But while you might decide not to play negative free bids after a $1\dz$ or 1M opening, after a Polish $1\cz$ it is very important to have a way to get 5-card suits into the auction, and we use either CFBs (after most overcalls) or transfers (after $1\sz$) to achieve this.

A CFB shows a hand which is prepared to play in that contract when opener has the weak option of his $1\cz$ opening -- that is, we will have a hand which would take out into that suit if partner opened a weak NT. The suit must certainly be at least 5 cards, but the bid does not generally promise much in the way of suit quality. For example, after $1\cz$ : ($2\dz$) we would be happy to bid $2\sz$ with $\sz$QTxxx $\hz$KQx $\dz$xx $\cz$xxx. We are expecting at least a doubleton in support if partner has the weak option.\footnote{1=4=3=5 shape is possible in theory, but those hands come up so rarely that we should not let this possibility affect responder's bidding in competition.} Occasionally a 5-2 fit may play poorly, but making a competitive bid and getting our suit into the auction is well worth the risk.

We have to be slightly more cautious if our suit is diamonds, because then there is a fair chance that partner will have shortage even in the weak option. So, with a poor 5-card diamond suit, we would look for alternative ways to bid the hand. Having said that, it would be wrong to wait until we pick up a 6-card suit before making the bid: the hands where partner has diamond shortage are not that common, and he could still have four (even five) diamonds in the weak option, where we would lose out if unable to find our fit.

When the CFB is a jump, it promises a hand which is happy for opener to compete to the 3-level on most hands with 3-card support. If the suit is not good enough for this, or the hand is too defensive, we would bid the suit at the 1-level instead.

The range for a CFB is about 5--10 HCP. It must be less than a game-forcing hand, because opener will pass most of the time with the weak type. On the other hand, the CFB is not a complete sign-off opposite the weak type, as opener is allowed to raise holding good support, and so it is still possible to reach game when all responder needs is a decent minimum with a good fit. So it is not necessary to stretch to make a stronger bid on too many 9- or 10-point hands.

The lower limit of 5 HCP comes from the fact that opener will force to game opposite a CFB with the strong option of $1\cz$ (indeed, any positive action from responder is game-forcing opposite the strong option). Not all 5-counts will be good enough, but if we have a good suit it is worth making a CFB in order to get the suit into the auction even if this would mean reaching a few light games opposite misfitting strong hands.

\vspace{4mm}
Continuations after a CFB are essentially the same as the methods that we use after fourth-seat interference over our positive responses. Even if opener's RHO actually passes, the situation is similar to over an overcall because opener is not forced to bid. However, if opener's RHO does make a bid we do not use support doubles after a CFB -- responder is already promising 5 cards, so if opener has 3-card support and a good enough hand he can normally make a natural raise. Instead, a double of an overcall after a CFB is take-out, which as usual tends to mean a fairly balanced hand without a stop. Similalrly, if a CFB is doubled, opener's redouble shows a strong hand (interested in taking a penalty): it is not a support redouble.

Apart from doubles, all the methods over fourth-seat interference apply; in particular:

\begin{itemize}
\item A bid of $3\cz$ is natural and not forcing.
\item A bid of $3\dz$ is artificial (unless it is a jump, or diamonds is partner's suit): if another cue-bid is available then $3\dz$ shows a game force with clubs;\footnote{Note that a $3\dz$ bid may show clubs even if that is the opponents' suit (so long as a $3\cz$ cue is available). For example, this would apply in the auction $1\cz$ : ($2\cz$) : $2\sz$ : (Pass) , $3\dz$. This is for practical reasons: our $1\cz$ opening is confusing to many players, and so it seems best to have a way to show a good hand with clubs in case they did not intend their bid as natural. We will not get any help from the TD if opponents have misbid, so we need to be able to find our own contracts.} otherwise $3\dz$ is a general game-forcing bid.
\item Bids in new major suits at the lowest available level are natural and game-forcing.
\item A bid of 3NT is to play. 2NT is also natural and forcing; if partner's suit is a minor then coninuations after 2NT are artificial ($3\cz$ waiting, $3\dz$/3M show shortage).
\item After a major-suit CFB, a $4\dz$ rebid shows a serious slam try in the major.
\item After a $2\cz$ CFB, opener's $2\dz$ rebid shows club support, as it would in the uncontested sequence $1\cz$ : $2\cz$ , $2\dz$. Similarly, after a $2\dz$ CFB, opener's $2\hz$  rebid is artificial and shows the natural option of $1\cz$.
\end{itemize}

\subsection{Responder's take-out double}

At the 1-level, responder's take-out double shows about the same strength as a normal positive response.

\begin{description}
\item[1$\cz$ : (1$\dz$) : Dbl] shows both majors.
\item[1$\cz$ : (1$\hz$) : Dbl] \emph{denies} 4 spades.\footnote{Some systems use the take-out double of $1\hz$ to distinguish between hands with 4 spades and hands with 5+ spades. This is not necessary after a Polish $1\cz$ because we can bid a competitive $2\sz$ on hands with 5 spades. So our $1\sz$ bid will tend to show exactly 4 spades unless strong, and we can use the double for something else.}
\item[1$\cz$ : (1$\sz$) : Dbl] will usually have exactly 4 hearts.
\end{description}
Because we can show hands with 5+ hearts with a transfer over $1\sz$, there is usually no need to double on such hands. So a double of $1\sz$ will only very rarely have 5 hearts.

At the 2-level, only slightly more strength is needed for a take-out double, provided that the hand has suitable shape. However, the partnership will not be able to stop below game if opener has the natural variant. Unlike at the 1-level, we may need to double on some hands with a 5-card major -- when our suit is higher than the opponents' suit, one way to force to game is to double and then bid our suit at the 3-level. In particular we may have to do this when we have a fairly balanced game-forcing hand but do not have a stop in the opponents' suit (if we do have a stop then we will be able to transfer to our suit at the 3-level and then bid 3NT). The take-out double at the 2-level also includes nearly all game-forcing balanced hands with 4 cards in an unbid major (or both majors).

At the 3-level or higher, a take-out double requires close to an opening hand. Here we will very often need to double on game-forcing hands with a 5-card suit, because there is no other way to make a forcing bid with a suit.

\vspace{4mm}
In replying to a take-out double, any suit bid from opener at the minimum level shows the weak type. Even after an intervening bid, opener with the weak hand may still bid a suit provided that this is no more than one level higher than would have been possible without the intervention.

A jump response to the take-out double tends to show the natural option of $1\cz$. With the strong option, opener will usually have to start with a cue-bid, although some other bids such as a jump in NT clearly show the strong type.

Special meanings are given to opener's `responsive' double if opener's RHO raises his partner's suit:
\begin{itemize}
\item If the suit is raised to the 2-level, for example $1\cz$ : ($1\sz$) : Dbl : ($2\sz$), opener's double artificially shows the natural type. (We do not bid $3\cz$ with this hand, because opener will very often want to compete in clubs with the weak type.)
\item If the suit is raised to the 3-level, for example $1\cz$ : ($1\sz$) : Dbl : ($3\sz$) or $1\cz$ : ($2\dz$) : Dbl : ($3\dz$), opener's double is value-showing and effectively replaces the cue-bid which has been taken away.
\end{itemize}

We also need to consider opener's minimum rebid in no-trumps. If opener's RHO raised the level of the auction then this shows a strong hand, but otherwise the bid can be weak. In response to a 1-level take-out double, a 1NT bid shows the weak type and is natural. Also in response to a 3-level take-out double, opener can bid 3NT on the weak type, but this bid is more wide-ranging and not particularly well defined, being used on a variety of stronger hands as well. Most interesting is the 2NT response to a 2-level take-out double: this is forcing, and is used on any of the following hands:

\begin{itemize}
\item (Only over a double of 2M) A `scrambling' weak hand, offering a choice of contracts at the 3-level if responder is not game-forcing.
\item A minimum weak hand wanting to play in the suit below the one doubled.
\item A hand with 5+ clubs and a stopper in opponents' suit, wanting to keep a club slam in the picture (will rebid 3NT next, or $4\cz$ with a huge hand).
\item A strong balanced hand with a 4-card major suit (will cue-bid next). This implies a stop in the opponents' suit, and offers a choice between 3NT and 4M.
\end{itemize}
So if opener bid the suit below the one doubled naturally at the 3-level, this shows a more constructive hand than going through 2NT, although still the weak type.

Note that over a double of 2M, if opener has a strong balanced hand with 4 cards in the other major but without a stop in the opponents' suit, he should be prepared to bid to 4M and not worry about offering the choice of playing in 3NT. Responder would not normally double with a stop in the opponents' suit unless he had 4 cards in the other major as well.

\vspace{4mm}
If responder makes a low-level double and follows up by bidding a new suit on the next round, this is non-forcing if it is below 3 of the suit doubled. A bid at the 3-level would be invitational -- if responder wants to bid a suit competitively he might be able to use 2NT Lebensohl. At lower levels the bid of a new suit may be just competing, but `correcting' partner's response is more constructive.

If the doubler bids a new suit \emph{above} 3 of the suit doubled, this is forcing to game and promises 5 cards. Responder would only bid this way on a hand which would be difficult to describe starting with a transfer: normally fairly balanced but without a stop in the opponents' suit.

\subsection{After a double of $1\cz$}

When they double our $1\cz$ opening we make a number of changes to our usual system:
\begin{itemize}
\item All 2-of-a-suit bids are CFBs, and 3-of-a-suit is pre-emptive. We use the \emph{redouble} to take care of strong hands with clubs: it shows 5+ clubs and at least invitational values. Further bidding after the redouble is natural except that opener's $2\dz$ rebid (or $3\dz$ if opponents take $2\dz$ away) shows a good hand with club support, like after a $2\cz$ CFB.
\item The 1NT response is natural, but with a range of about 7-11 HCP, not the narrow invitational range that we have in an uncontested sequence. Now opener's rebid of $2\cz$ shows the natural option of $1\cz$ but is not forcing; $2\dz$ is artificial and shows a game force with clubs. If opener's RHO makes a bid then we use the normal system over fourth-seat interference, so for example a non-jump $3\cz$ bid is non-forcing, with the $3\dz$ convention showing a game force with clubs.
\item A pass shows 3+ clubs, and normally a negative hand. Opener's rebids of $1\hz$ or above are the same as after a normal negative, for example opener's $2\dz$ is an artificial game force. Opener can also pass the double, or bid $1\dz$ on a minimum with 4 or 5 diamonds, or redouble for rescue.
\end{itemize}
The $1\dz$ and 1M responses are mostly unchanged by the double, but there are some hands which have been taken out: negatives with 3+ clubs and some balanced hands are taken out of $1\dz$, and some hands with a 5-card suit may choose to make a CFB rather than bidding $1\dz$ or 1M. However, this would not affect the meaning of subsequent bids -- we play `system on'.


\subsection{After a 1-level overcall}

Like after a double, responder's 1NT bid in competition is natural and shows about 7-11 HCP. Continuations are the same as for $1\cz$ : (Dbl) : 1NT, with $2\cz$ being non-forcing and $2\dz$ showing a game-forcing hand with clubs.

Most of responder's other methods have been described already. Bidding 1M is natural, with continuations the same as in an uncontested sequence. CFBs apply at the 2-level after a $1\dz$ or $1\hz$ overcall. But over a $1\sz$ overcall, we play transfer-based methods instead:

\begin{description}
\item[2$\cz$] shows 5+ diamonds.
\item[2$\dz$] shows 5+ hearts.
\item[2$\hz$] shows 5+ hearts, invitational.
\item[2$\sz$] shows 5+ clubs and at least invitational values.
\end{description}
With a weaker hand with clubs we bid $3\cz$, which is effectively a 3-level CFB (and over this, $3\dz$ is game forcing with club support, as usual). Of course, this requires a better suit than a 2-level CFB, but a 5-card suit may still be sufficient if no other call looks attractive.

The minimum requirements for the $2\cz$ and $2\dz$ transfers are exactly the same as for a normal CFB. Opener's rebids are the same as for the corresponding CFBs as well,\footnote{Note in particular that $1\cz$ : ($1\sz$) : $2\cz$ : (Pass/Dbl) , $2\hz$ is artificial showing clubs.} with the obvious exception that if opener's RHO passes and opener has a hand which would pass a CFB, he has to complete the transfer. Also, because the transfers are not limited, opener has to be slightly more careful when bidding game, so as not to pre-empt any slam investigation.

The main difference is when responder gets to make his rebid. If opener completes the transfer and responder rebids in a new suit, this is forcing. Raising his own suit is invitational.

The invitational $2\hz$ bid is also very similar to a CFB.

The transfer to clubs works slightly differently. Here opener usually bids 2NT with the weak type, and now responder will bid $3\cz$ with an invite; opener's $3\cz$ also shows the weak type but with genuine clubs. Other bids show strong hands, but some strong hands can also start with 2NT in order to keep the bidding low and see what happens. A jump to 3NT shows a strong hand but will be a relatively poor one for slam.

\vspace{4mm}
The cue-bid to show an invitational or better hand with 5+ clubs is not just played over a $1\sz$ overcall: it also applies over $1\dz$ or $1\hz$. Again, opener will bid 2NT or $3\cz$ with the weak type; other bids show stronger hands, including 2M.

We also need a way to show good hands with diamonds after a $1\hz$ overcall. Here we play

\begin{description}
\item[3$\cz$] shows a game-forcing hand with 5+ diamonds.
\item[3$\dz$] shows an invitational hand with 6+ diamonds.
\end{description}
After either of these, opener's $3\sz$ rebid would be natural and strong.

So, by using transfers for suits lower than the overcall, we always have a way to make a forcing bid with a 5-card suit. Without a 5-card suit, we can bid 2NT (natural and forcing to game) or make a jump cue-bid (asking for a stop, almost certainly balanced with no 4-card major), or double if the shape is right.

New suits at the 3-level, if not already defined, are natural and pre-emptive.

\begin{center}
\begin{tabular}{|l|p{3cm}|p{3cm}|p{3cm}|} \hline
 & \st Over $1\cz$ : ($1\dz$) & Over $1\cz$ : ($1\hz$) & Over $1\cz$ : ($1\sz$) \\ \hline
Dbl & \st both majors & denies 4 spades & usually 4 hearts  \\ \hline
$1\hz$ & \st 4+ $\hz$s, forcing & & \\ \hline
$1\sz$ & \st 4+ $\sz$s, forcing & 4+ $\sz$s, forcing & \\ \hline
1NT & \st natural & natural & natural \\ \hline
$2\cz$ & \st CFB & CFB & 5+ $\dz$s \\ \hline
$2\dz$ & \st 5+ $\cz$s, inv+ & CFB & 5+ $\hz$s, weak or GF \\ \hline
$2\hz$ & \st CFB & 5+ $\cz$s, inv+ & 5+ $\hz$s, invitational \\ \hline
$2\sz$ & \st CFB & CFB & 5+ $\cz$s, inv+ \\ \hline
2NT & \st natural GF & natural GF & natural GF \\ \hline
$3\cz$ & \st pre-emptive & 5+ $\dz$s, GF & CFB \\ \hline
$3\dz$ & \st GF asking for stop & 5+ $\dz$s, invitational & pre-emptive \\ \hline
$3\hz$ & \st pre-emptive & GF asking for stop & pre-emptive \\ \hline
$3\sz$ & \st pre-emptive & pre-emptive & GF asking for stop \\ \hline
3NT & \st to play & to play & to play \\ \hline
\end{tabular}
\end{center}

Methods over a 1NT overcall are rather simpler: double is for penalties, 2-level bids are CFBs and 3-level bids are pre-emptive. 2NT shows a game-forcing two-suiter. This is all the same as in most standard systems.

\subsection{After a 2-level overcall}

Here, in addition to the usual take-out double and 2-level CFBs, we play transfers starting at 2NT.

These methods are very similar to those that might be used for dealing with interference over a 1NT opening. This is no coincidence, because our $1\cz$ opening will usually be a weak NT (particularly if responder has a good enough hand to take a bid) and this is what responder will play for until proven otherwise. Indeed it would make a lot of sense to take whatever methods we normally use after a 1NT opening and apply those here. However, transfer-based methods are clearly superior and so we define them as part of our standard $1\cz$ system. Best is to adopt these methods over $1\cz$ and then use them over 1NT as well.

The rules for transfers are as follows:
\begin{itemize}
\item When responder transfers to a suit lower than the opponents' suit, opener must complete the transfer with the weak type. Responder will then pass if he has a purely competitive hand; bidding on is game-forcing.
\item The transfer to the opponents' suit is, of course, not needed as a genuine transfer, and we redefine it to be a natural invitational bid.
\item The transfer to a suit higher than the opponents' suit promises at least invitational values. So opener may bid game with a maximum weak hand. To show the strong type opener must make some other bid, for example a cue to show support. 4NT here is a natural bid -- we would start with a cue if we wanted to bid Blackwood.
\item A $3\sz$ bid (`transfer to 3NT') asks for a stop in the opponents' suit.
\end{itemize}
So the complete table of bids starting from 2NT is:

\begin{center}
\begin{tabular}{|l|p{2cm}|p{2cm}|p{2cm}|p{2cm}|} \hline
 & \st $1\cz$ : ($2\cz$) & $1\cz$ : ($2\dz$) & $1\cz$ : ($2\hz$) & $1\cz$ : ($2\sz$) \\ \hline
2NT & \st nat. inv. & 5+ $\cz$s & 5+ $\cz$s & 5+ $\cz$s \\ \hline
$3\cz$ & \st 5+ $\dz$s, inv+ & 5+ $\cz$s, inv. &  5+ $\dz$s & 5+ $\dz$s \\ \hline
$3\dz$ & \st 5+ $\hz$s, inv+ & 5+ $\hz$s, inv+ & 5+ $\dz$s, inv. & 5+ $\hz$s \\ \hline
$3\hz$ & \st 5+ $\sz$s, inv+ & 5+ $\sz$s, inv+ & 5+ $\sz$s, inv+ & 5+ $\hz$s, inv. \\ \hline
$3\sz$ & \st stop ask & stop ask & stop ask & stop ask \\ \hline
3NT & \st to play & to play & to play & to play \\ \hline
\end{tabular}
\end{center}

\vspace{3mm}
After a transfer to a suit higher than the opponents' suit, responder has no way to ask for a stop after the transfer is completed. So, with a fairly balanced game-forcing hand and no stop in the opponents' suit, he may have to start with a double intending to rebid his 5-card suit on the next round.

\subsection{After a 3-level or higher overcall}

Responder's suit bids are all non-forcing here: we would have to double with a strong hand or simply bid a game. Generally bidding is the same as if the opponents had opened a pre-emptive bid, except that we have slightly more information about opener's hand.

After responder's non-forcing suit bids, opener's jump to 4NT is natural. To set responder's suit as trumps with a strong hand, opener would cue-bid.

\subsection{Responder's rebid after a negative}

If opener has shown natural clubs or a strong type then it should be fairly clear what responder's calls mean. We need to be more concerned about what happens when opener could still have the weak type, as in auctions like these:

\begin{itemize}
\item $1\cz$ : (Pass) : $1\dz$ : ($1\sz$) , Pass : (Pass) : ?
\item $1\cz$ : (Pass) : $1\dz$ : ($1\hz$) , Pass : ($2\hz$) : ?
\item $1\cz$ : (Pass) : $1\dz$ : (Pass) , $1\hz$ : ($2\cz$) : ?
\end{itemize}

Responder may always bid a new suit at the 1- or 2-level with the genuine negative hand. So for example $1\cz$ : (Pass) : $1\dz$ : ($1\hz$) , Pass : ($2\hz$) : $2\sz$ shows at most 5 HCP with 5+ good spades, probably short in hearts. However, reponder is not allowed to bid a natural suit at the 3-level with a negative. Bids at the 3-level will always have diamonds, unless that is the opponents' suit:
\begin{description}
\item[3$\cz$] If this is not a jump, it shows 7-10 HCP with 5 diamonds and 4+ clubs. If it is a jump then it is a game-forcing minor two-suiter.
\item[3$\dz$] This is always natural and invitational.
\item[3M] Forcing to game, implying 5+ diamonds and 4 cards in the major.
\item[cue-bid] Forcing to game, promising 5 diamonds, usually no stop in opponents' suit.
\end{description}

However, in the protective position, responder can bid 2NT Lebensohl over a 2M overcall. This gives responder a way to take action on a negative hand when holding a suit lower than the opponents' suit. This does not apply when responder's RHO made a bid: now 2NT shows a good hand with diamonds including a stop in the opponents' suit. If we hold this hand in the protective position we can start with 2NT and follow up with a strong rebid.

Responder's double of a suit bid is for take-out. At the 1- or 2-levels, we might double on any of the following hands:
\begin{itemize}
\item Any game-forcing balanced hand, or fairly balanced hand with diamonds.
\item A balanced hand in the 7-10 HCP range with doubleton in the opponents' suit.
\item In the protective position only, a negative hand with shortage in the opponents' suit.
\end{itemize}
After a double of 2M, opener's 2NT is scrambling. Opener should not be expecting 4 cards in an unbid major, since apart from the negative, these hands would have bid 1M. So the take-out double is primarily focused on the minor suits.

All of the game-forcing balanced hands not wanting to declare no-trumps start with a double. They can later cue-bid to ask for a stop, or bid 3 of an unbid major artificially to show a stop in opponents' suit and ask partner to bid 3NT.

At the 3-level or higher, we may have to double with either the balanced hand or the game-forcing hand with diamonds.

\subsection{Responder's rebid in other situations}

As in the previous section, we shall only look at situations where opener could still have the weak type.

First of all, suppose that the initial response was $1\hz$ or $1\sz$. In a competitive situation after a 1M response, responder's new suit bid is not forcing, unless it is above 3 of his first suit. Minor-suit bids normally show 5 cards there. For bids in new suits at the 3-level:
\begin{itemize}
\item If competing over an opponent's $2\sz$ bid (when our suit is hearts), a 3m bid is purely competitive with probably 5-5 shape.
\item Otherwise, a 3m bid is invitational, with 5 or 6 cards in the minor and normally only 4 in the major.
\end{itemize}
Responder's 2NT rebid over an opponent's $2\dz$/$2\hz$/$2\sz$ is Lebensohl. Whether in protective position or not, this shows a weakish hand wanting to compete in a minor. He is likely to have only 4 cards in his first suit, particularly when opponents' interference is $2\sz$ where we would bid 3m with a 5-5 shape. Responder can also double for take-out, with either 4 or 5 cards in his major. After the take-out double of 2M, opener's 2NT is scrambling.

2NT Lebensohl also applies if responder's first bid was a transfer, for example $1\cz$ : ($1\sz$) : $2\dz$ : ($2\sz$) , Pass : (Pass) : 2NT  shows a hand which wants to compete in a minor suit. A new suit bid at the 3-level would be \emph{forcing}.

Take-out doubles apply almost everywhere. For example, responder's double in the sequence $1\cz$ : ($1\hz$) : Pass : ($2\hz$) , Pass : (Pass) : Dbl is for take-out, implying a hand too weak to take action over $1\hz$. Provided that we are in the protective position, we can double at the 2-level even with a negative. A suit bid from responder after passing initially must also be a very weak hand, since we would have made a CFB or a transfer with positive values.

\subsection{When responder is a passed hand}
The big difference when responder is a passed hand is that all transfers are off. So after a $1\sz$ overcall, for example, responder's bids of $2\cz$, $2\dz$ and $2\hz$ become CFBs. However the use of a cue-bid to show a good hand with clubs still applies (and responder can still just about have a game force in this case).

A natural game-forcing 2NT bid also does not really make sense as a passed hand. This could be given a different meaning, depending on what our pre-emptive openings would have meant.

\newpage
\section{Dealing with artificial interference}

Because our $1\cz$ opening is artificial, in most countries the opponents are allowed to play any defence they like to it. The Polish $1\cz$ opening is nothing like a Strong Club, since most of the time it is a weak NT, and so opponents would be well advised to concentrate on constructive methods. However many opponents like to wheel out whatever crazy defence to a Strong Club they are playing, or perhaps just want to make full use of the lack of system restrictions. We should be happy for them to do this, because a destructive defence should cause them more problems than it causes us, but in order to take advantage we have to be properly prepared.

One type of interference is easy to handle: if they \emph{double}, our bids do not depend on what the meaning of their double is. Even if their double of $1\cz$ shows the majors (say), we can still bid major suits naturally.\footnote{Sometimes the opponents won't be sure of what their double means anyway. It might be best not to ask.} Also if opponents bid no-trumps, it rarely makes a difference to us what their bid means -- our double shows values and responder's new suits are non-forcing. So for the rest of this section, we shall be assuming that we are dealing with an artificial \emph{suit bid}.

\subsection{Artificial $1\dz$ and $1\hz$ overcalls}

When opponents make an artificial $1\dz$ or $1\hz$ overcall, we use double to show the suit that they bid.

Apart from the change to the meaning of a double, responder's bids are normally `system on' as if their overcall was natural. Even if, for example, their $1\dz$ overcall shows diamonds and hearts, we can still bid $1\hz$ naturally. And we still have artificial sequences to show minor-suit hands: playing `system on' means that --
\begin{itemize}
\item $1\cz$ : ($1\dz$) : $2\dz$ or $1\cz$ : ($1\hz$) : $2\hz$ shows an invitational or better hand with 5+ clubs.
\item $1\cz$ : ($1\hz$) : $3\cz$ shows a game-forcing hand with 5+ diamonds.
\end{itemize}
So normally we do not have a way to show a CFB in the suit bid by the opponents. However, if their overcall shows a specific suit, then we change things round so that the cue-bid of that suit shows the good hand with clubs, and 2 of the suit that they bid becomes a CFB.

\subsection{Artificial $1\sz$ overcalls}

When a $1\sz$ overcall is artificial, this does \emph{not} change the meaning of a double: it still take-out, normally with precisely 4 hearts. If the doubler later bids spades, this is always an artificial `cue-bid'.

The usual system of transfers still applies, with one exception: the $2\hz$ bid becomes a transfer to spades. So, over an artificial $1\sz$ overcall:
\begin{description}
\item[2$\cz$] Shows 5+ diamonds.
\item[2$\dz$] Shows 5+ hearts.
\item[2$\hz$] Shows 5+ spades.
\item[2$\sz$] Shows 5+ clubs, invitational or better.
\end{description}

Sometimes the opponents' $1\sz$ bid shows length in a specific suit. This does not affect responder's initial bid, and he can transfer to the suit to show length there. But if responder makes some other bid and then later bids the opponents' suit, this is an artificial cue.

In general, opponents who play artificial overcalls at the 1-level tend not to require particularly good suits for their bids. After a bid showing `spades and clubs', for example, it is still perfectly possible that we might want to play in spades ourselves. There are also some very nebulous bids such as a $1\sz$ showing `3 or 4 spades' which we would treat as not showing any particular suit at all.

\subsection{Artificial overcalls at the 2-level or higher}

The meaning of a double here depends on what exactly their overcall shows:
\begin{itemize}
\item If the overcall shows two specific suits, not including the suit bid, then double is value-showing.
\item If the overcall shows two specific suits including the suit bid, then double is take-out.
\item If the overcall promises one specific suit (not the suit bid), then double is take-out of the suit shown.
\item If the overcall does not promise any specific suit, then double is take-out of the suit \emph{bid}. (But if their overcall is a multi-way bid and none of the options includes the suit bid, the double will tend to look more like a value-showing double.)
\end{itemize}
In the last of these cases, it is possible that opener might still have length in the suit doubled, but responder will not have. As with any take-out double, if the doubler subsequently rebids in the suit doubled, this is an artificial cue. For example, if the opponenets overcall our $1\cz$ with a $2\hz$ bid showing `either hearts or the black suits', then responder can double as take-out of hearts; if he later rebids $3\hz$ this shows a hand which suspects that the opponents do have hearts and needs to ask for a stopper there.

Apart from the double, responder's calls mostly follow the same rules as for natural interference. If the interference was at the 2-level, then transfers start at 2NT and go up to $3\sz$. The transfer to the suit bid is now a `real' transfer -- as with the transfers to lower suits, opener must complete the transfer with the weak type, and now any further bid from responder is forcing to game. However, if the opponents' bid shows a particular suit or suits:
\begin{itemize}
\item The transfer to a suit that opponents have shown normally changes its meaning to a natural invitational bid. (This is the same rule that is used for natural interference.) But if this suit has \emph{also} been shown by the opponents (for example $3\hz$ after $1\cz$ : $2\cz$ showing the majors) then this `transfer' shows a stop in the suit you are bidding, whereas $3\sz$ would imply a stop in the other suit.
\item If we can cue-bid opponents' suit at the 2-level, this shows 5+ clubs, at least invitational (i.e.~the same as a 2-level cue after natural 1-level interference).
\item If opponents have shown \emph{two} suits that we can cue-bid at the 2-level, then the lower one shows 5+ clubs, whereas the higher one shows the other unbid suit, again at least invitational.
\end{itemize}

\subsection{Opener's rebid}

If partner has bid no-trumps or shown a suit then there is usually no problem. Our methods over interference do not rely on opponents having shown a suit; in some cases the $3\dz$ convention may take the place of a cue-bid. Note that if partner bid 1M naturally, then opener's double of a 1- or 2-level bid is always a support double, even if their bid was artificial. This includes the double of a cue-bid.

If partner passed over an overcall, or bid the negative:
\begin{itemize}
\item Opener's reopening actions always treat the overcall as if it was natural. So after $1\cz$ : ($1\sz$) : Pass : (Pass), even if the $1\sz$ bid showed a random hand, we would double as take-out of spades and cue-bid to show clubs, as if they had made a natural spade bid. However we can show a game-forcing hand with spades by starting with $2\dz$.
\item If RHO makes a suit bid that promises length in a different suit, then double is take-out of the suit shown, and a cue-bid shows clubs.
\item If RHO makes a suit bid that does not promise length in any particular suit, then double is take-out of the suit bid. For example if LHO bids $2\dz$ to show `either hearts or the black suits' and RHO bids $2\hz$ `pass or correct', then opener's double is take-out of hearts. A double followed by a cue-bid of that suit on the next round is always artificial. However opener can make a direct cue-bid (e.g. bidding $3\hz$ over $2\hz$ here) to show a natural hand.
\end{itemize}
Alternatively, if partner made a double of an artificial bid:
\begin{itemize}
\item If RHO passes, then we would of course pass with length in the suit bid. So a cue-bid is artificial and shows a strong hand.
\item If RHO bids a new suit, whatever this means, we play double as penalties, and a `cue-bid' of the suit RHO bid is artificial. A bid of LHO's suit is natural and shows the strong option of $1\cz$.
\item If RHO makes a natural `raise' of the suit LHO bid (for example, bidding $2\sz$ to play over a random $1\sz$ bid), then we play the same system as if the overcall had been natural, so double shows clubs. However, if RHO's raise is artificial (e.g.~pass or correct), then we revert to the system for dealing with artificial interference, using double for penalties.
\end{itemize}

\subsection{`Undiscussed' club bids}

We will often be faced with a situation where opponents are unsure what their bids mean, most commonly with the $2\cz$ overcall of $1\cz$. Recommended is to treat this in the same way as we would treat an artificial $2\cz$ overcall meaning `clubs or not clubs'. That is, most of our bids are the same as if $2\cz$ was natural, with double being take-out of clubs, but responder can show clubs by bidding a 2NT transfer.

If responder makes a double or CFB over $2\cz$, we will continue as if the club bid was natural. So opener's $3\cz$ rebid is a cue. But if opener has the $3\dz$ convention available, for example in a sequence such as $1\cz$ : ($2\cz$) : $2\sz$ : (Pass) , $3\dz$, this shows a game force with clubs. (We defined this sequence as showing clubs even when their $2\cz$ bid was supposedly natural.)

Similarly, if opener's RHO makes an `undiscussed' $2\cz$ bid, it is best to treat a double as take-out of clubs (or support, if partner bid 1M), but have opener's $3\cz$ bid as natural.

\newpage
\section{Optional methods}

\subsection{Transfer checkback after $1\cz$ : 1M , 1NT}

This is a method designed to give responder as many descriptive sequences as possible. It is designed to work with a style where $1\cz$ : $1\hz$ , 1NT could have 4 spades: the transfer methods give us ways to find our spade fits.

The methods include a transfer to the major that responder has already bid. Obviously this has no effect on who is declarer, as it would do over a 1NT opening: the reason we use this transfer is to give us more bidding sequences.

\begin{center}
\begin{tabular}{|l|p{5cm}|p{5cm}|} \hline
 & \st Over $1\cz$ : $1\hz$ , 1NT & Over $1\cz$ : $1\sz$ , 1NT \\ \hline
$2\cz$ & \st Forces $2\dz$ unless partner is minimum with 3-card support & Forces $2\dz$ unless partner is minimum with 3-card support \\ \hline
$2\dz$ & \st Shows 5+ hearts & Shows 5+ spades and 4+ hearts (initially asking for preference) \\ \hline
$2\hz$ & \st 5 hearts and 4 spades, weak, asking for preference & Shows 5+ spades \\ \hline
$2\sz$ & \st Precisely 4-4 majors, invitational & To play \\ \hline
2NT & \st Transfer to clubs & Transfer to clubs \\ \hline
3m & \st Natural and invitational & Natural and invitational \\ \hline
\end{tabular}
\end{center}

\vspace{3mm}
Note that the sequence $1\cz$ : $1\hz$ , 1NT : $2\hz$, while still being somewhat like a transfer in that it shows the suit above, is \emph{not} forcing. The other transfers are forcing; after opener completes the transfer, any rebid from responder is forcing to game (or possibly an invitational hand with a 6-card major might transfer and then rebid 3M). The transfers are not used on invitational hands with a 5-card suit, since these go through $2\cz$. This works best if the partnership's normal system over a 1NT opening involves the 2NT rebid after a transfer being game-forcing (as in `Keri' or `re-transfers') since then we can play system on here.

The $2\cz$ bid will be made on one of the following types of hands:
\begin{itemize}
\item Weak hands with 5+ diamonds, intending to pass $2\dz$.
\item Invitional hands. After opener bids $2\dz$, responder's rebid of $2\hz$, $2\sz$, 2NT, $3\cz$ or $3\dz$ is invitational, with bids in new suits promising at least 5 cards in the original major. (With only 4 cards in the major, we would bid the second suit directly over 1NT).
\item Game-forcing balanced hands with precisely 4 cards in the major, not suitable for bidding 3NT directly. This could be because responder is 4-4 in the majors (in which case he rebids $3\sz$) or because the hand is worth a slam try (in which case he can rebid a 4-card minor at the 4-level).
\end{itemize}
Opener usually bids $2\dz$ over $2\cz$, but if he has a minimum with 3 cards in partner's major he bids 2M instead. This ensures that we play in the right place when responder is 5-5, and sometimes the knowledge that partner has 3-card support is all responder needs to bid game.

Note that after $1\cz$ : $1\sz$ , 1NT responder has two ways to get out in $2\sz$. We would bid a direct $2\sz$ when we want to play there even opposite a 1=4=3=5 hand. Whereas, if responder bids $2\hz$ transfer, opener will bid 2NT rather than completing the transfer on the 1=4=3=5 shape, and now responder's 3m bids are to play. This also gives responder valuable information when he has the game-forcing type.

\subsection{Transfers over a $3\cz$ overcall}

Generally we believe that competitive free bids are more useful than natural forcing free bids over 3-level interference. However, we could use transfers to give us a way to show both a competitive hand and a forcing hand. This works best over $3\cz$, where all we are giving up is the competitive bid in diamonds:

\begin{description}
\item[3$\dz$] shows hearts.
\item[3$\hz$] shows spades.
\item[3$\sz$] shows diamonds.
\end{description}

A similar thing would work when their overcall is $3\dz$, but now we would lose the ability to bid a non-forcing $3\hz$.

\subsection{Opener's $4\dz$ response to a take-out double}

It is often quite difficult for opener to show a strong hand in response to partner's take-out double. Perhaps the worst auction is $1\cz$ : ($3\sz$) : Dbl, where we might want to bid $4\hz$ on a minimum, but at the same time we have no way to show a good hand with hearts without going past $4\hz$.

A possible solution is to use $4\dz$ to show a `good' $4\hz$ bid. Opener is unlikely to want to bid a natural $4\dz$ showing a minimum: most balanced hands will pass or perhaps bid 3NT instead, and hands with 5 diamonds are rare -- if opener has a hand which wants to bid a minor, it is much more likely that the suit is clubs.

Even when opener does have a cue-bid available, or when a double can be used to show a strong hand as in the auction $1\cz$ : ($2\sz$) : Dbl : ($3\sz$) , Dbl, it would still make things clearer if $4\dz$ was used to show a good 4M bid. The natural meaning of $4\dz$ is of very little use here.

So it makes sense to play that whenever responder doubles a major suit, opener's bid of $4\dz$ shows a hand with length in the other major too strong for a natural 4M bid.

\end{document}